
\begin{exercise}[Sublimation ou fusion ?]

    \begin{questions}
        \question Comment peut-on expliquer que la fonte de la banquise ou celle d'un glacier d'altitude conduise à de l'eau liquide, alors qu'un cirrus (nuage de haute altitude formé de cristaux de glace) disparaisse au soleil sans donner de pluie ?
            
        \question Le dioxyde de carbone possède un point triple de coordonnées $(T_T, P_T) = (-57~ \si{\celsius}, 5 \cdot 10^5 ~ \si{\pascal})$, quel changement d'état observe-t-on lorsqu'on laisse de la carboglace (\ce{CO2} solide) à pression ambiante ?
            
    \end{questions}

\end{exercise}